% D1_COMMUTATION
\begin{theorem}[Normalized analytic--convex commutation theorem]
\label{thm:r3-d1-commutation}
Assume (H1)--(H4).  Let \(A\Subset\mathbb R^k\) be a compact target set
contained in the interior of the constraint mean image
\[
 \{CDQ(\Theta+C^*\lambda):
 \Theta\in B',\ \lambda\in V\}
\]
for smaller source and multiplier balls, with the saddle staying a positive
distance from their boundaries.  Then:
\begin{enumerate}[label=(\roman*)]
\item \(Q_\varepsilon\to Q\) locally uniformly with all Fr\'echet
      derivatives on every smaller complex ball.
\item The good rate is
\[
 I(x)=\sup_{\Theta\in X}
 \{\langle\Theta,x\rangle-Q(\Theta)+Q(0)\}.
\]
\item For \(a\in A\),
\[
 Q^a(\Theta)
 =\inf_\lambda\{Q(\Theta+C^*\lambda)-\lambda\cdot a\}
 -\inf_\lambda\{Q(C^*\lambda)-\lambda\cdot a\}
\]
has a unique analytic source-dependent minimizer \(\lambda_\Theta^a\).
\item Its Hessian is
\[
 D^2Q^a
 =Q_{\Theta\Theta}
 -Q_{\Theta\lambda}Q_{\lambda\lambda}^{-1}
 Q_{\lambda\Theta}.
\]
\item For a typed continuous or exponentially approximable contraction
      \(T:X^*\to E\),
\[
 Q_T(\eta)=Q(T^*\eta),
 \qquad I_T(y)=\inf_{Tx=y}I(x),
\]
and derivatives commute with \(T\).
\item Constraint and contraction commute with the necessary normalization:
\[
 \boxed{
 ((Q^a\circ T^*)^*)(y)
 =\inf_{Tx=y,\ Cx=a}I(x)
 -\inf_{Cx=a}I(x).}
\]
The identity extends from exposed regular points by lower-semicontinuous
closure.
\end{enumerate}
\end{theorem}

\begin{proof}
Normal holomorphic convergence on a larger ball and Cauchy's formula prove
(i).  The full LDP gives the Laplace upper bound and the lower bound at exposed
points; the rate-dense exposed approximation in (H2) proves (ii).

For \(a\in A\), the critical equation is
\[
 CDQ(\Theta+C^*\lambda_\Theta^a)=a.
\]
Uniform positivity of \(CD^2QC^*\), the interior-image assumption, and the
boundary margin give existence and uniqueness throughout the declared compact
target set.  The analytic inverse-function theorem gives (iii), and
implicit differentiation plus the envelope theorem gives (iv).

Part (v) follows from the exact finite-volume identity
\(\langle\eta,T\mathcal X_\varepsilon\rangle
 =\langle T^*\eta,\mathcal X_\varepsilon\rangle\), the contraction theorem,
and the chain rule.  For (vi), the normalized constrained rate is
\[
 I^a(x)=
 \begin{cases}
 I(x)-I_C(a),&Cx=a,\\
 +\infty,&Cx\ne a,
 \end{cases}
 \qquad I_C(a)=\inf_{Cx=a}I(x).
\]
Contracting \(I^a\) under \(T\) gives the boxed formula.  This constant is
essential whenever \(a\) is atypical.
\end{proof}

% D1_LIKELIHOOD
\subsection{Finite-mean Gaussian tangents and exact local likelihoods}

Fix a regular source \(\Theta\) and put
\[
 m_{\varepsilon,\Theta}=DQ_\varepsilon(\Theta),
 \qquad
 m_\Theta=DQ(\Theta),
 \qquad
 \Sigma_\Theta=D^2Q(\Theta).
\]
Define the exactly centered field
\[
 Z_\varepsilon
 =\sqrt{\mu_\varepsilon}
 (\mathcal X_\varepsilon-m_{\varepsilon,\Theta}).
\]

\begin{theorem}[Gaussian tangent and exact likelihood-ratio convergence]
\label{thm:r3-d1-likelihood}
Every fixed finite-dimensional projection of \(Z_\varepsilon\) converges to a
centered Gaussian vector with covariance \(\Sigma_\Theta\).  For a local
source direction \(h\), define
\[
 \begin{aligned}
 L_\varepsilon(h)=\exp\Big\{&
 \sqrt{\mu_\varepsilon}
 \langle h,\mathcal X_\varepsilon-m_{\varepsilon,\Theta}\rangle\\
 &-\mu_\varepsilon\big[
 Q_\varepsilon(\Theta+h/\sqrt{\mu_\varepsilon})
 -Q_\varepsilon(\Theta)
 -\mu_\varepsilon^{-1/2}DQ_\varepsilon(\Theta)h
 \big]\Big\}.
 \end{aligned}
\]
Then \(L_\varepsilon(h)\) is the exact Radon--Nikodym density of the
\(\Theta+h/\sqrt{\mu_\varepsilon}\) law relative to the \(\Theta\) law and
has expectation one for every \(\varepsilon\).  Jointly,
\[
 (Z_\varepsilon,L_\varepsilon(h))
 \Rightarrow
 \left(Z,
 e^{\langle h,Z\rangle-rac12\Sigma_\Theta(h,h)}
 \right).
\]
The likelihoods are uniformly integrable and the tilted Gaussian mean is
\(\Sigma_\Theta h\).

If, in addition, the finite fields satisfy the B3 nuclear-space tightness and
the C2 resolved-kernel convergence hypotheses, the same conclusion holds in
the declared path topology and for the likelihood martingale process.
\end{theorem}

\begin{proof}
Under the \(\Theta\)-tilted law, the characteristic function of a finite
projection is
\[
 \exp\left\{
 \mu_\varepsilon\left[
 Q_\varepsilon(\Theta+i t/\sqrt{\mu_\varepsilon})
 -Q_\varepsilon(\Theta)
 -i\mu_\varepsilon^{-1/2}DQ_\varepsilon(\Theta)t
 \right]
ight\}.
\]
The linear term cancels exactly at finite volume.  Cauchy bounds on a larger
complex ball make the cubic remainder
\(O(\mu_\varepsilon^{-1/2}|t|^3)\), and convergence of the Hessian gives the
Gaussian characteristic function.

Expanding the exact real source density jointly with imaginary test directions
gives the likelihood limit.  Since the displayed exponent is algebraically
the difference between the two finite log densities, its expectation is one
without any rate assumption on
\(m_{\varepsilon,\Theta}-m_\Theta\).  A slightly larger real source ball gives
a uniform \(L^{1+\eta}\) bound.  Le Cam's third lemma gives the mean shift.
B3 tightness and C2 kernel convergence upgrade finite-dimensional convergence
to the stated process result.
\end{proof}

\begin{corollary}[Conditioned and contracted Gaussian tangent]
\label{cor:r3-d1-conditioned}
For \(a\in A\), the conditioned covariance is
\[
 \Sigma^a
 =\Sigma-\Sigma C^*(C\Sigma C^*)^{-1}C\Sigma.
\]
Every typed projection \(T\) has covariance \(T\Sigma^aT^*\), and its local
likelihood ratios are exact finite-volume Cameron--Martin approximants centered
at the corresponding finite conditioned mean.
\end{corollary}

\begin{proof}
Apply the Schur-complement part of
Theorem~\ref{thm:r3-d1-commutation} to the finite constrained pressure and then
Theorem~\ref{thm:r3-d1-likelihood}; apply the chain rule for \(T\).
\end{proof}
